\documentclass[11pt,a4paper]{article}
\usepackage[utf8]{inputenc}
\usepackage[T1]{fontenc}
\usepackage{lmodern}
\usepackage{amsmath,amssymb,amsfonts,amsthm}
\usepackage{geometry}
\geometry{margin=1in}
\usepackage{xcolor}
\usepackage{hyperref}
\usepackage{cleveref}
\usepackage{booktabs}
\usepackage{array}
\usepackage{cite}
\usepackage{microtype}
\usepackage{tcolorbox}
\usepackage{slashed}

\hypersetup{
    colorlinks=true,
    linkcolor=blue,
    citecolor=red,
    urlcolor=cyan,
    pdftitle={Testing a Micro/Macro Dual-Scale Hypothesis on K3 x T2: Scale-Factor Duality, the Cohen-Kaplan-Nelson Bound, and the Dark-Energy Length in Lean 4},
    pdfauthor={Xavier Callens and the SocrateAI Scientific Agora Collaboration}
}

\newtheorem{theorem}{Theorem}[section]
\newtheorem{lemma}[theorem]{Lemma}
\newtheorem{definition}[theorem]{Definition}
\newtheorem{proposition}[theorem]{Proposition}
\newtheorem{corollary}[theorem]{Corollary}
\newtheorem{conjecture}[theorem]{Conjecture}

\theoremstyle{remark}
\newtheorem{remark}[theorem]{Remark}

% ---------------------------------------------------------------------------
% Epistemic tier badges, following the convention of papers 7-8
% (papers/publication/paper{7,8}_*.tex) and docs/PAPER_REVISION_BRIEF_2026_09_17.md §0.
% ---------------------------------------------------------------------------
\newcommand{\tierA}{\colorbox{green!15}{\textsf{\scriptsize Tier A --- Lean kernel}}}
\newcommand{\tierL}{\colorbox{blue!10}{\textsf{\scriptsize Tier L --- literature}}}
\newcommand{\tierC}{\colorbox{orange!15}{\textsf{\scriptsize Tier C --- conjecture}}}

\newtcolorbox{scopebox}{%
  colback=gray!5, colframe=gray!45, boxrule=0.4pt, arc=1pt,
  left=6pt, right=6pt, top=3pt, bottom=3pt,
  fonttitle=\bfseries\small, title={Scope of Lean 4 verification}
}
\newtcolorbox{casebox}[1]{%
  colback=red!4, colframe=red!45!black, boxrule=0.4pt, arc=1pt,
  left=6pt, right=6pt, top=3pt, bottom=3pt,
  fonttitle=\bfseries\small, title={#1}
}

\title{\Huge \textbf{Testing a Micro/Macro Dual-Scale Hypothesis}\\[0.3em]
\LARGE \textbf{on $K3\times T^2$:}\\[0.2em]
Scale-Factor Duality, the Cohen--Kaplan--Nelson Bound,\\
and the Dark-Energy Length in Lean~4}

\author{\textbf{Xavier Callens}$^{1}$ \and \textbf{SocrateAI Scientific Agora Collaboration}$^{1}$ \\[0.5em]
$^{1}$\textit{SocrateAI Research Laboratory for Mathematical Physics \& Formal Verification}}

\date{September 2026}

% [layout] allow long Lean identifiers to break after underscores
\makeatletter\let\LMoldunderscore\_\renewcommand{\_}{\LMoldunderscore\allowbreak}\makeatother
\setlength{\emergencystretch}{3em}
% [/layout]
\begin{document}

\maketitle

\begin{abstract}
We report \texttt{DualScaleCosmology} (Lean~4, toolchain \texttt{v4.33.1}, Mathlib), a companion
library that formally tests one reading of a ``micro/macro dual-scale'' hypothesis motivating this
research programme: that a microscopic length $\ell_{\mathrm{micro}}\sim\ell_P$ and a macroscopic
length $\ell_{\mathrm{macro}}\sim H_0^{-1}$ are linked by non-perturbative duality. The hypothesis
is not, and cannot be, a theorem; we instead formalize the two pieces of the literature that
actually bear on it --- string-cosmology scale-factor duality (Gasperini--Veneziano) and the
Cohen--Kaplan--Nelson (CKN) ultraviolet/infrared bound --- and test their numerical and structural
consequences. The results are mixed and mostly negative for the literal hypothesis. Taken
literally, $\ell_P$ as the ultraviolet cutoff at $H_0^{-1}$ fails the CKN bound by more than
$10^{30}$ (Tier~A arithmetic on Tier~L constants). Read instead as a T-dual pair, the two scales
saturate the CKN bound exactly at their self-dual length $s=\sqrt{\ell_P\,c/H_0}\approx47\,\mu$m;
we then prove, as an exact identity, that $s$ equals the observed dark-energy length
$\Lambda^{-1/4}$ up to a fixed, computable factor $(8\pi/3\Omega_\Lambda)^{1/4}$ --- so the
numerical coincidence often quoted between such lengths is algebra, not independent evidence. As a
string (Regge) scale $s$ is excluded by collider data by more than $10^{30}$ in $\alpha'$; as the
radius of a single extra dimension it is disfavored only by an order-one factor, indistinguishable
at this level of rigor from the ``dark dimension'' proposal of Montero--Vafa--Valenzuela. We also
report, as a matter of methodological record, a citation error caught only by computing a numeric
consequence of it, and an infrastructure misconfiguration that silently starved a local prover of
its models for most of a work session. All Lean content is Tier~A (44 declarations, 31 theorems,
0 \texttt{sorry}, standard axioms only); every physical identification is stated as Tier~L
(literature) or Tier~C (this project's own reading), never conflated with the former. We close with
a concrete list of open problems and proposed research directions for the community.
\end{abstract}

\newpage
\tableofcontents
\newpage

%%%%%%%%%%%%%%%%%%%%%%%%%%%%%%%%%%%%%%%%%%%%%%%%%%%%%%%%%%%%%%%%%%%%%%%%%%%%%%%
\section{Introduction}
\label{sec:intro}
%%%%%%%%%%%%%%%%%%%%%%%%%%%%%%%%%%%%%%%%%%%%%%%%%%%%%%%%%%%%%%%%%%%%%%%%%%%%%%%

\subsection{The hypothesis under study}
Perturbative superstring theories fix the fundamental string tension $T=1/(2\pi\alpha')$ near the
Planck scale, $\ell_s\sim\ell_P\sim10^{-35}$~m. A dual-scale framework, motivating this
programme's earlier papers~\cite{Callens2026Paper1,Callens2026Paper7}, proposes instead a pair of
dynamically coupled length scales: a microscopic scale $\ell_{\mathrm{micro}}\sim\ell_P$
(quantum-gravity corrections, ultraviolet completion) and a macroscopic scale
$\ell_{\mathrm{macro}}\sim H_0^{-1}$ or an intermediate galactic/dark-sector scale, linked by
non-perturbative duality transformations, scale-invariance breaking, or topological defects. This
is \emph{not} a theorem anywhere in the literature, and this paper does not treat it as one. It is
a Tier~C proposal, and the question this paper answers is narrower and more disciplined: which
pieces of existing physics literature actually relate a micro and a macro length by a precise
mathematical statement, and what do their kernel-checked consequences say about the hypothesis as
stated?

\subsection{Two literature anchors, not a re-derivation}
Two, and — after a systematic search — only two, published results were found that relate a
Planck-scale ultraviolet input to a cosmological infrared scale with a precise statement rather
than a general expectation.

\begin{enumerate}
\item \textbf{Scale-factor duality.} Gasperini and Veneziano~\cite{GasperiniVeneziano1993} show
that the scale factor $a(t)$ of an FRW background is odd under $a\mapsto a^{-1}$ in exactly the
$O(d,d)$ sense that a compactification radius is odd under Buscher T-duality $R\mapsto\alpha'/R$
(paper~1's Theorem~2.1). This is a genuine duality, but it is a symmetry of the early-universe
string effective action, not on its own a statement relating today's Hubble scale to the Planck
scale.
\item \textbf{The Cohen--Kaplan--Nelson bound.} Cohen, Kaplan, and Nelson~\cite{CKN1999} show that
an effective field theory with ultraviolet cutoff $\Lambda$ in a box of infrared size $L$, demanding
no state already collapsed to a black hole, satisfies $L^3\Lambda^4\lesssim LM_P^2$. This is an
actual inequality --- not a duality --- tying a micro (UV) scale to a macro (IR) scale, and the
authors themselves instantiate it with $L$ at the current cosmological horizon. If any published
result deserves the name ``micro/macro dual scale,'' this is it, and it predates this project's
hypothesis by 28~years.
\end{enumerate}
Nothing else in the literature search (cosmic-string tension bounds, the swampland distance
conjecture, the ``dark dimension'' proposal) relates a Planck-scale and a Hubble-scale length by a
precise, checkable statement rather than an order-of-magnitude expectation; those sources are used
below, but as tests of the two anchors above, not as further anchors.

\subsection{Contents and how to read the tier labels}
\Cref{sec:tiers} restates this project's epistemic tier system. \Cref{sec:duality} formalizes
scale-factor duality. \Cref{sec:ckn} formalizes the CKN bound, including a citation error caught
and corrected in the course of this work. \Cref{sec:falsify} tests the literal hypothesis
numerically and finds it fails. \Cref{sec:selfdual} finds where the two anchors meet: a self-dual
length at which the CKN bound is saturated. \Cref{sec:readings} gives that length two physical
readings --- excluded as a string scale, disfavored (not excluded) as an extra-dimension radius ---
and \Cref{sec:identity} proves it is, as an exact identity, the observed dark-energy length.
\Cref{sec:strings} formalizes cosmic F-strings, the one genuinely micro object that can be
cosmologically long. \Cref{sec:towers} tests whether the swampland distance conjecture is an
independent third route, or shares scale-factor duality's structure. \Cref{sec:outofscope} records
what was deliberately \emph{not} formalized, and why. \Cref{sec:scope} summarizes what is and is
not proved. \Cref{sec:methods} reports two verification-discipline case studies. \Cref{sec:open}
and \Cref{sec:directions} close with open problems and proposed research directions for the
community. \Cref{app:repro} is a reproducibility appendix.

%%%%%%%%%%%%%%%%%%%%%%%%%%%%%%%%%%%%%%%%%%%%%%%%%%%%%%%%%%%%%%%%%%%%%%%%%%%%%%%
\section{Epistemic tiers}
\label{sec:tiers}
%%%%%%%%%%%%%%%%%%%%%%%%%%%%%%%%%%%%%%%%%%%%%%%%%%%%%%%%%%%%%%%%%%%%%%%%%%%%%%%
As in this programme's earlier papers, every claim below carries one of three tags.
\begin{itemize}
  \item \tierA\ --- accepted by the Lean~4 kernel: \texttt{\#print axioms} reports only
  \texttt{propext}, \texttt{Classical.choice}, \texttt{Quot.sound}; no \texttt{sorry}, no
  \texttt{admit}, no \texttt{native\_decide}, no project-declared \texttt{axiom}.
  \item \tierL\ --- established in the cited literature, quoted with a file and line range in
  \texttt{papers/foundations/*.txt} that a reader can \texttt{grep}, not re-derived here.
  \item \tierC\ --- a modeling choice or conjecture introduced by this project. \textbf{The
  identification of a Lean object with a physical quantity is never Tier~A}, no matter how the
  surrounding prose reads.
\end{itemize}
Two conventions specific to this paper, both forced by incidents recorded in
\Cref{sec:methods}: every numeric bound quoted from a source is rounded in the direction that
\emph{weakens} the conclusion being drawn (never in the direction that helps it); and every
Lean theorem name is checked to be plain ASCII, because this project's own audit tooling silently
mis-parses non-ASCII identifiers (\Cref{sec:methods}).

%%%%%%%%%%%%%%%%%%%%%%%%%%%%%%%%%%%%%%%%%%%%%%%%%%%%%%%%%%%%%%%%%%%%%%%%%%%%%%%
\section{Scale-factor duality}
\label{sec:duality}
%%%%%%%%%%%%%%%%%%%%%%%%%%%%%%%%%%%%%%%%%%%%%%%%%%%%%%%%%%%%%%%%%%%%%%%%%%%%%%%
Gasperini--Veneziano~\cite{GasperiniVeneziano1993} (ll.~372--373): ``the Hubble parameter is odd
under scale factor duality, $H\to-H$, so that a duality transformation $a\to a^{-1}$ maps an
expanding solution'' to a contracting one; the self-dual case is ``a metric that satisfies
$a^{-1}(t)=a(-t)$'' (ll.~639--640). \tierL.

\begin{theorem}[Scale-factor duality is an involution] \label{thm:invol} \tierA\
For $a\ne0$, \textnormal{(}\texttt{ScaleFactorDuality.scaleFactorDual\_invol}\textnormal{)}
$(a^{-1})^{-1}=a$.
\end{theorem}

\begin{theorem}[Self-dual point] \label{thm:fixed} \tierA\
For $a>0$, $a^{-1}=a\iff a=1$ \textnormal{(}\texttt{scaleFactorDual\_fixed\_iff}\textnormal{)} ---
the $t=0$ instance of $a^{-1}(t)=a(-t)$.
\end{theorem}

\begin{theorem}[The Hubble parameter is odd under the duality] \label{thm:hubble} \tierA\
For $a:\mathbb{R}\to\mathbb{R}$ differentiable at $t$ with $a(t)\ne0$, with
$H(t):=a'(t)/a(t)$,
\[
H\big(s\mapsto a(s)^{-1}\big)(t)\;=\;-\,H(a)(t) \qquad \textnormal{(\texttt{hubble\_dual})}.
\]
\end{theorem}
This is a genuine differentiation, via Mathlib's \texttt{HasDerivAt.inv}, not an arithmetic
shadow. An earlier draft of this work left \Cref{thm:hubble} as Tier~L, on the assumption that
real differentiation was out of scope for this project's usual style; that assumption was checked
against confirmed proof work in a sibling repository and found wrong (\Cref{sec:methods}).

\begin{theorem}[Transferred dual-scale bound] \label{thm:sum} \tierA\
For $a>0$, $a+a^{-1}\ge2$ \textnormal{(}\texttt{cosmoDualScale\_ge\_two}\textnormal{)}, the exact
algebraic transfer of Stream~2's torus bound
$\operatorname{tr}G+\operatorname{tr}G^{-1}\ge2d$~\cite{Callens2026Paper8} to the scale factor.
\end{theorem}

%%%%%%%%%%%%%%%%%%%%%%%%%%%%%%%%%%%%%%%%%%%%%%%%%%%%%%%%%%%%%%%%%%%%%%%%%%%%%%%
\section{The Cohen--Kaplan--Nelson bound}
\label{sec:ckn}
%%%%%%%%%%%%%%%%%%%%%%%%%%%%%%%%%%%%%%%%%%%%%%%%%%%%%%%%%%%%%%%%%%%%%%%%%%%%%%%

\begin{casebox}{Case study: a citation error caught only by computing something}
An earlier commit of this work stated CKN's eq.~(2) as $L^3\Lambda^4\le M_P^2$. This compiled as
valid Lean and was reported, correctly, as Tier~A --- but it is \emph{dimensionally
inconsistent} ($[L^3\Lambda^4]=\mathrm{GeV}$, $[M_P^2]=\mathrm{GeV}^2$) and therefore not what the
paper says. The kernel cannot catch a false Tier~L citation; only someone re-reading the source
can. The error was found not by re-reading in the abstract, but by trying to \emph{use} the
statement for the numeric test of \Cref{sec:falsify} and finding the units did not close.
Re-extracting the source with \texttt{pdftotext~-layout} (rather than trusting the first,
unlayouted pass) gives the actual eq.~(2): $L^3\Lambda^4\lesssim L\,M_P^2$
(ll.~77--78)~\cite{CKN1999}, dimensionally consistent, and — dividing by $L>0$ — exactly the
well-known holographic-dark-energy scaling $L^2\Lambda^4\lesssim M_P^2$. Fixing the exponent
\emph{simplified} the Lean proofs below: no fractional-power (\texttt{Real.rpow}) machinery is
needed for the corrected statement, only one \texttt{nlinarith} step.
\end{casebox}

CKN, ll.~77--78, 80~\cite{CKN1999}: $L^3\Lambda^4\lesssim L\,M_P^2$, so ``the IR cutoff scales
like $\Lambda^{-2}$.'' \tierL.

\begin{theorem}[CKN bound, simplified form] \label{thm:cknbound} \tierA\
For $L>0$, $L^3\Lambda^4\le LM^2\implies L^2\Lambda^4\le M^2$ \textnormal{(}\texttt{CKNBound.ckn\_bound}\textnormal{)}.
\end{theorem}

\begin{theorem}[Solved for the combination $L\Lambda^2$] \label{thm:cknla} \tierA\
For $L,\Lambda,M>0$, $L^2\Lambda^4\le M^2\implies L\Lambda^2\le M$
\textnormal{(}\texttt{ckn\_L\_Lambda\_sq\_le}\textnormal{)}: since $L^2\Lambda^4=(L\Lambda^2)^2$, this
is a one-step square-root extraction, not a fractional power.
\end{theorem}

\begin{corollary}[Solved for each variable] \label{cor:cknvars} \tierA\
$L\le M/\Lambda^2$ \textnormal{(}\texttt{ckn\_L\_le}\textnormal{)} and $\Lambda^2\le M/L$
\textnormal{(}\texttt{ckn\_Lambda\_sq\_le}\textnormal{)}.
\end{corollary}

%%%%%%%%%%%%%%%%%%%%%%%%%%%%%%%%%%%%%%%%%%%%%%%%%%%%%%%%%%%%%%%%%%%%%%%%%%%%%%%
\section{Testing the literal hypothesis}
\label{sec:falsify}
%%%%%%%%%%%%%%%%%%%%%%%%%%%%%%%%%%%%%%%%%%%%%%%%%%%%%%%%%%%%%%%%%%%%%%%%%%%%%%%
CKN's own worked example, ll.~113--114~\cite{CKN1999}: fixing the IR cutoff $L$ at the current
cosmological horizon, eq.~(2) implies $\Lambda\sim10^{-2.5}\,\mathrm{eV}$. \tierL. We use this
number directly, rather than re-deriving $H_0^{-1}$ in Planck units from CODATA constants, to
avoid re-introducing exactly the kind of unit-conversion error of the case study above.

\begin{theorem}[The literal hypothesis fails by $\ge10^{30}$] \label{thm:falsify} \tierA\
With $\Lambda_{\mathrm{CKN}}:=3.2\times10^{-3}\,\mathrm{eV}$ (CKN's value, rounded \emph{up},
conservative) and $M_P:=1.2\times10^{28}\,\mathrm{eV}$ (CODATA 2018 Planck energy, rounded
\emph{down}, conservative),
\[
M_P\;\ge\;10^{30}\,\Lambda_{\mathrm{CKN}} \qquad\textnormal{(\texttt{CKNInstance.planckEnergy\_gt\_ckn\_horizon\_cutoff})}.
\]
\end{theorem}
The true ratio is $\approx3.9\times10^{30}$; both roundings understate it. \textbf{Reading:} if the
Stream~3 hypothesis's ``$\ell_{\mathrm{micro}}\sim\ell_P$'' is the ultraviolet cutoff the CKN bound
would allow at ``$\ell_{\mathrm{macro}}\sim H_0^{-1}$,'' it is not --- by 30 orders of magnitude.
This is a direct, literature-sourced counter-example to the literal reading of the hypothesis.

%%%%%%%%%%%%%%%%%%%%%%%%%%%%%%%%%%%%%%%%%%%%%%%%%%%%%%%%%%%%%%%%%%%%%%%%%%%%%%%
\section{Where the bound meets the duality: the self-dual length}
\label{sec:selfdual}
%%%%%%%%%%%%%%%%%%%%%%%%%%%%%%%%%%%%%%%%%%%%%%%%%%%%%%%%%%%%%%%%%%%%%%%%%%%%%%%
\Cref{thm:falsify} tests one specific reading. This section asks what \emph{does} follow if
$\ell_{\mathrm{micro}}$ and $\ell_{\mathrm{macro}}$ are instead read as a T-dual pair --- the
reading \Cref{sec:duality} actually formalizes --- with $\alpha':=\ell_{\mathrm{micro}}\ell_{\mathrm{macro}}$
and CKN's Planck mass identified with $M=1/\ell_{\mathrm{micro}}$ (natural units). Both
identifications are this project's choices, stated here as explicit hypotheses, not results.
\tierC.

\begin{theorem}[The pair is a scale-factor-duality pair] \label{thm:pairisdual} \tierA\
In units of $s:=\sqrt{\ell_{\mathrm{micro}}\ell_{\mathrm{macro}}}$,
$\ell_{\mathrm{macro}}/s=(\ell_{\mathrm{micro}}/s)^{-1}$
\textnormal{(}\texttt{SelfDualCutoff.selfDual\_pair\_is\_scaleFactorDual}\textnormal{)}.
\end{theorem}

\begin{theorem}[CKN $\iff$ UV length at least the self-dual length] \label{thm:ckniff} \tierA\
With $L=\ell_{\mathrm{macro}}$, $M=1/\ell_{\mathrm{micro}}$, $\Lambda=1/\ell_{\mathrm{UV}}$,
\[
L^2\Lambda^4\le M^2 \;\iff\; \ell_{\mathrm{micro}}\,\ell_{\mathrm{macro}}\le\ell_{\mathrm{UV}}^2
\qquad\textnormal{(\texttt{ckn\_iff\_uvLength\_ge\_selfDual})},
\]
saturated at equality \textnormal{(}\texttt{ckn\_saturated\_at\_selfDual}\textnormal{)}. Being an
\textsf{iff} between non-equivalent-looking statements, this cannot be vacuous: it fails in each
direction for suitable $\ell_{\mathrm{UV}}$.
\end{theorem}

\begin{proposition}[Numeric bracket] \label{prop:bracket} \tierA\
With $\ell_P=1.616255\times10^{-35}$~m (CODATA 2018) and $c/H_0=1.3672\times10^{26}$~m
($H_0=67.66$~km/s/Mpc, Planck~2018 VI~\cite{Planck2018VI}, Table~2, via
\texttt{astropy.cosmology.Planck18}),
\[
40\,\mu\mathrm{m}\;\le\;\sqrt{\ell_P\,c/H_0}\;\le\;70\,\mu\mathrm{m}
\qquad\textnormal{(\texttt{selfDual\_length\_sq\_bracket})},
\]
and CKN's own horizon-scale cutoff length $\hbar c/\Lambda_{\mathrm{CKN}}$ lies in the same
bracket \textnormal{(}\texttt{cknHorizon\_length\_bracket}\textnormal{)}: $s\approx47.0\,\mu$m,
$\hbar c/\Lambda_{\mathrm{CKN}}\approx62\,\mu$m.
\end{proposition}
\textbf{Reading:} the hypothesis's ``$\ell_{\mathrm{micro}}\sim\ell_P$'' is not where the CKN
bound's implied micro cutoff sits (\Cref{thm:falsify}); the geometric mean
$s=\sqrt{\ell_P\,\ell_{\mathrm{macro}}}$ is.

%%%%%%%%%%%%%%%%%%%%%%%%%%%%%%%%%%%%%%%%%%%%%%%%%%%%%%%%%%%%%%%%%%%%%%%%%%%%%%%
\section{Two physical readings of the self-dual length}
\label{sec:readings}
%%%%%%%%%%%%%%%%%%%%%%%%%%%%%%%%%%%%%%%%%%%%%%%%%%%%%%%%%%%%%%%%%%%%%%%%%%%%%%%

\subsection{As a string scale: excluded}
\label{sec:strings-scale}
If $\alpha'=s^2$ is the fundamental string's Regge slope, then, since $M^2=(N-1)/\alpha'$ at
level $N$~\cite{Tong2009} (§1.3), string excitations sit near $\hbar c/s\approx4$~meV. CMS-EXO-19-012,
137~fb$^{-1}$ at 13~TeV, ll.~1044--1048~\cite{CMS2020}: string resonances (Regge excitations of
quarks and gluons, l.~63) are excluded below 7.9~TeV at 95\%~CL, in the low-string-scale models of
its own references. \tierL.

\begin{theorem}[Excluded as a Regge slope] \label{thm:regge} \tierA\
$\ell_P\,c/H_0\;\ge\;10^{30}\,(\hbar c/7.9\,\mathrm{TeV})^2$ \textnormal{(both sides
$\mathrm{m}^2$; \texttt{CosmicString.selfDual\_alphaPrime\_exceeds\_cms\_ceiling}}\textnormal{)}.
\end{theorem}
This excludes only the \emph{Regge-slope} reading; it says nothing about $s$ as a
\emph{compactification radius} (\Cref{sec:strings-radius}), and it is model-dependent on CMS's own
low-string-scale models, stated as such.

\subsection{As an extra-dimension radius: disfavored, not excluded}
\label{sec:strings-radius}
Torsion-balance tests (Lee et al., ll.~285--289~\cite{LeeAdelberger2020}): gravitational-strength
Yukawa interactions are excluded for range $<38.6\,\mu$m; a single large extra dimension's toroidal
radius must be $<30\,\mu$m at 95\%~CL. \tierL.

\begin{theorem}[Above the torsion-balance bound] \label{thm:torsion} \tierA\
$(1.5\times30\,\mu\mathrm{m})^2\le\ell_P\,c/H_0$ \textnormal{(}\texttt{DarkEnergyScale.selfDual\_above\_torsion\_radius\_bound}\textnormal{)}.
\end{theorem}
So identifying $s$ directly with a large extra-dimension radius exceeds the experimental bound by a
factor of only $\approx1.6$ --- an order-one margin, not the $10^{30}$ of \Cref{thm:regge}. This
paper does \emph{not} claim the extra-dimension reading is excluded.

%%%%%%%%%%%%%%%%%%%%%%%%%%%%%%%%%%%%%%%%%%%%%%%%%%%%%%%%%%%%%%%%%%%%%%%%%%%%%%%
\section{The self-dual length is the dark-energy length}
\label{sec:identity}
%%%%%%%%%%%%%%%%%%%%%%%%%%%%%%%%%%%%%%%%%%%%%%%%%%%%%%%%%%%%%%%%%%%%%%%%%%%%%%%
With $\rho_\Lambda=\Omega_\Lambda\cdot3H_0^2/(8\pi G)$ (critical-density form used by
Copeland--Kibble, eq.~(4.1), ll.~430--434~\cite{CopelandKibble2009}) and $G=\ell_P^2$
($\hbar=c=1$):

\begin{theorem}[An exact identity] \label{thm:identity} \tierA\
\[
\rho_\Lambda^{-1}\;=\;\frac{8\pi}{3\Omega_\Lambda}\,(\ell_P\,H_0^{-1})^2
\qquad\textnormal{(\texttt{DarkEnergyScale.rhoLambda\_inv\_eq})},
\]
i.e.\ $\rho_\Lambda^{-1/4}=(8\pi/3\Omega_\Lambda)^{1/4}\,s$: the self-dual length \emph{is} the
dark-energy length $\Lambda_{\mathrm{cc}}^{-1/4}$, up to a fixed, $\ell_P$- and
$H_0$-independent number. Equivalently, in Montero--Vafa--Valenzuela's (MVV) parameterization
$l=\lambda\,\Lambda_{\mathrm{cc}}^{-1/4}$ (l.~67~\cite{MonteroVafaValenzuela2023}), $s$ has
$\lambda_s=(3\Omega_\Lambda/8\pi)^{1/4}$, \textnormal{\textbf{a pure number}}
\textnormal{(}\texttt{selfDual\_lambda4\_eq}\textnormal{)}.
\end{theorem}

\begin{casebox}{Why this blocks a natural-looking argument}
MVV quote $\Lambda_{\mathrm{cc}}^{-1/4}\approx88\,\mu$m (l.~303). CKN's own horizon cutoff length is
$\approx62\,\mu$m (\Cref{prop:bracket}), and $s\approx47.0\,\mu$m. \emph{These are not three
independent numbers agreeing}: by \Cref{thm:identity} they are the same formula evaluated with
different order-one conventions (reduced vs.\ non-reduced Planck mass gives a factor
$(8\pi)^{1/4}\approx2.2$ in length). A numeric bracket test alone
(\texttt{darkEnergyLength4\_bracket}, giving $85$--$90\,\mu$m for $\Lambda_{\mathrm{cc}}^{-1/4}$)
would sit \emph{outside} the earlier $[40,70]\,\mu$m bracket of \Cref{prop:bracket} --- which by
itself is evidence the bracket is convention-dependent, not physical. No document in this
programme should cite the numerical closeness of $s$, the CKN horizon length, and
$\Lambda_{\mathrm{cc}}^{-1/4}$ as independent corroboration.
\end{casebox}

\begin{proposition}[$\lambda_s$ vs.\ MVV's range] \label{prop:lambda} \tierA\
$\lambda_s^4\in[0.5^4,1)$, i.e.\ $\lambda_s\approx0.535$
\textnormal{(}\texttt{selfDual\_lambda\_above\_mvv\_range}\textnormal{)}: inside MVV's broad range
$10^{-4}<\lambda<1$ (l.~424) but $\ge5\times$ their central estimate's upper end $\lambda\sim10^{-1}$
(ll.~424--427). MVV separately exclude a light \emph{string} tower at this scale, on different
grounds (local EFT breakdown, ll.~313--316) --- consistent with \Cref{thm:regge}.
\end{proposition}

%%%%%%%%%%%%%%%%%%%%%%%%%%%%%%%%%%%%%%%%%%%%%%%%%%%%%%%%%%%%%%%%%%%%%%%%%%%%%%%
\section{Cosmic F-strings: a genuine micro object at macro scale}
\label{sec:strings}
%%%%%%%%%%%%%%%%%%%%%%%%%%%%%%%%%%%%%%%%%%%%%%%%%%%%%%%%%%%%%%%%%%%%%%%%%%%%%%%
A cosmic superstring is literally a micro object (string tension) stretched to cosmological
length --- the most direct candidate for the hypothesis's mechanism. Copeland--Myers--Polchinski
(CMP)~\cite{CopelandMyersPolchinski2004} give the ten-dimensional $(p,q)$ tension, eq.~(3.3),
ll.~502--506, and the four-dimensional tension $\mu=e^{2A}\bar\mu$ with warp factor $e^{2A}$,
eq.~(3.2), l.~497. \tierL. For the fundamental string, $(p,q)=(1,0)$, unwarped ($e^{2A}=1$, a
Tier~C simplification that can only \emph{lower} $\mu$ if relaxed), $G\mu/c^2=\ell_P^2/(2\pi\alpha')$.

\begin{theorem}[An observational bound on $G\mu$ is a lower bound on $\alpha'$] \label{thm:gmu} \tierA\
For $\alpha'>0$, $G\mu/c^2\le g\iff\ell_P^2\le2\pi g\,\alpha'$
\textnormal{(}\texttt{CosmicString.fStringGmu\_le\_iff}\textnormal{)}. For the self-dual pair,
$\alpha'=\ell_P L\implies G\mu/c^2=\ell_P/(2\pi L)$
\textnormal{(}\texttt{fStringGmu\_selfDual}\textnormal{)}.
\end{theorem}

\begin{proposition}[The self-dual $\alpha'$ gives a badly under-tense string] \label{prop:cmpwindow} \tierA\
$G\mu/c^2\le10^{-61}$ for the self-dual pair
\textnormal{(}\texttt{selfDual\_Gmu\_below\_cmp\_window}\textnormal{)}, more than 50 orders below
the lower end $10^{-11}$ of CMP's brane-inflation window (eq.~(5.1), ll.~782--783), whose own
caveat (ll.~809--811) notes the window assumes inflaton-driven perturbations.
\end{proposition}
\textbf{Not formalized}: the defect-formation and horizon-scaling mechanism (Kibble, network
scaling) that would let a cosmic string carry the macro scale in the first place ---
\Cref{sec:outofscope}.

%%%%%%%%%%%%%%%%%%%%%%%%%%%%%%%%%%%%%%%%%%%%%%%%%%%%%%%%%%%%%%%%%%%%%%%%%%%%%%%
\section{Dual towers and the swampland distance conjecture}
\label{sec:towers}
%%%%%%%%%%%%%%%%%%%%%%%%%%%%%%%%%%%%%%%%%%%%%%%%%%%%%%%%%%%%%%%%%%%%%%%%%%%%%%%
The Swampland Distance Conjecture (SDC, due to Ooguri--Vafa, quoted at
ll.~2167--2181~\cite{PaltiReview2019}): from any point $P$ in moduli space there is a point $Q$ at
infinite geodesic distance and an infinite tower of states with $M(Q)\sim M(P)e^{-\alpha
d(P,Q)}$, $\alpha>0$. \tierL.

\begin{casebox}{Why eq.~(2.100) itself is not formalized}
The SDC's own equation is asymptotic, with an unspecified $\alpha$ and an ``$\sim$''. Any Lean
theorem about the bare expression $M_0e^{-\alpha d}$ (positive, monotone, tends to zero) is a fact
about \texttt{Real.exp}, not physics, and would restate its own definition --- the exact vacuous-check
failure a sibling project in this ecosystem retracted after finding it in its own work (see
\Cref{sec:methods}). We formalize instead Palti's \emph{motivating remark} just above the SDC box
(ll.~2158--2161): dual towers whose mass product is constant, under an explicit Tier~C model
$M_1M_2=M_0^2$.
\end{casebox}

\begin{theorem}[The dual-tower structure is P3.1's involution] \label{thm:towersum} \tierA\
Under $M_1M_2=M_0^2$ (Tier~C), $M_1+M_2\ge2M_0$
\textnormal{(}\texttt{DualTower.dualTower\_sum\_ge}\textnormal{)}, proved by rescaling to
$a=M_1/M_0$ and invoking \Cref{thm:sum} \emph{unchanged}. So the dual-tower structure is not an
independent piece of algebra: it is scale-factor duality again.
\end{theorem}

\begin{proposition}[Negative controls] \label{prop:negctrl} \tierA\
Without $M_1M_2=M_0^2$, both $\forall M_0,M_1,M_2{>}0,\,2M_0\le M_1{+}M_2$ and
$\forall M_0{>}0,\,\min(M_1,M_2)\le M_0$ are \emph{false}, by explicit kernel-checked
counterexample \textnormal{(}\texttt{dualTower\_sum\_ge\_needs\_product},
\texttt{dualTower\_min\_le\_needs\_product}\textnormal{)}. \Cref{thm:towersum} is therefore not
vacuous.
\end{proposition}

\begin{proposition}[Bridge to the exponential form] \label{prop:exptower} \tierA\
$(M_0e^{-\alpha d})(M_0e^{\alpha d})=M_0^2$ \textnormal{(}\texttt{expTowerPair\_product}\textnormal{)}
--- a fact about \texttt{Real.exp}, connecting the SDC's shape to the model above, not a
formalization of the SDC itself.
\end{proposition}

%%%%%%%%%%%%%%%%%%%%%%%%%%%%%%%%%%%%%%%%%%%%%%%%%%%%%%%%%%%%%%%%%%%%%%%%%%%%%%%
\section{Closed as out of scope: defect networks and domain walls}
\label{sec:outofscope}
%%%%%%%%%%%%%%%%%%%%%%%%%%%%%%%%%%%%%%%%%%%%%%%%%%%%%%%%%%%%%%%%%%%%%%%%%%%%%%%
How a defect network grows to carry the macro scale is a \emph{simulation} result, not a theorem:
string networks reach a scaling regime $\xi/t\approx\text{const}$ by ``general agreement from
several simulations'' (Copeland--Kibble, ll.~421--423~\cite{CopelandKibble2009}), with
$\rho_{\mathrm{str}}/\rho=8\pi G\mu\gamma^2/(3\nu^2)$ (eq.~(4.1)); domain walls reach
$\rho_{\mathrm{wall}}=A\sigma/t$ with $A\approx0.8\pm0.1$ from field-theory simulations
(Saikawa, eqs.~(2.17)--(2.18), ll.~269--283~\cite{Saikawa2017}). \tierL. Formalizing
$\rho_{\mathrm{str}}=\mu/\xi^2$ or $\rho_{\mathrm{wall}}=A\sigma/t$ in Lean would only restate
their own defining equations --- exactly the vacuity failure this project's own gate exists to
catch (\Cref{sec:methods}) --- so neither is formalized. The one structural fact recorded, Tier~L,
not proved: a defect network carries the macro scale \emph{by causality}
($\xi\lesssim t$~\cite{CopelandKibble2009}, ll.~272--276), while the micro scale enters only
through the tension; the two are not related by a duality in this mechanism, unlike
\Cref{sec:duality}.

%%%%%%%%%%%%%%%%%%%%%%%%%%%%%%%%%%%%%%%%%%%%%%%%%%%%%%%%%%%%%%%%%%%%%%%%%%%%%%%
\section{What is, and is not, proved}
\label{sec:scope}
%%%%%%%%%%%%%%%%%%%%%%%%%%%%%%%%%%%%%%%%%%%%%%%%%%%%%%%%%%%%%%%%%%%%%%%%%%%%%%%
\begin{scopebox}
\textbf{Is proved (Tier~A, kernel-checked, 44 declarations / 31 theorems, 0 \texttt{sorry},
standard axioms only):} that scale-factor duality is an involution with self-dual point $a=1$ and
odd Hubble parameter; that the corrected CKN bound $L^2\Lambda^4\le M^2$ follows from CKN's own
eq.~(2) and implies $L\Lambda^2\le M$; the numeric gap $\ge10^{30}$ between CKN's own horizon
cutoff and the Planck energy; that the CKN bound holds \emph{iff} the UV length is at least a
T-dual pair's self-dual length, with that length bracketed at $40$--$70\,\mu$m; that this length
is excluded as a Regge slope by $\ge10^{30}$ in $\alpha'$ and lies only $\approx1.6\times$ above a
torsion-balance bound as an extra-dimension radius; the \emph{exact identity} between that length
and the dark-energy length; a cosmic F-string's tension formula and its numeric position 50+ orders
below a literature brane-inflation window; that the dual-tower structure behind the swampland
distance conjecture is scale-factor duality's own involution, with kernel-checked negative
controls.

\textbf{Is not proved, anywhere in this paper:} the micro/macro dual-scale hypothesis itself, in
any reading; that any of the Tier~C identifications used above (T-dual pair, $M=1/\ell_{\mathrm{micro}}$,
$M_1M_2=M_0^2$) is physics rather than a modeling choice; the SDC's own asymptotic statement,
eq.~(2.100); the defect-network scaling mechanism or domain-wall dynamics (\Cref{sec:outofscope});
that a string scale of $\sim50\,\mu$m or an extra dimension of that size exists; and, as always in
this project, that any Lean object \emph{is} the physical system it is named after.
\end{scopebox}

%%%%%%%%%%%%%%%%%%%%%%%%%%%%%%%%%%%%%%%%%%%%%%%%%%%%%%%%%%%%%%%%%%%%%%%%%%%%%%%
\section{Methods: two verification-discipline case studies}
\label{sec:methods}
%%%%%%%%%%%%%%%%%%%%%%%%%%%%%%%%%%%%%%%%%%%%%%%%%%%%%%%%%%%%%%%%%%%%%%%%%%%%%%%
Beyond the citation-transcription error of \Cref{sec:ckn}, two further incidents are recorded here
because they generalize beyond this paper.

\subsection{A false ``models still resident'' claim}
An early version of this work's workflow log stated that both local prover models were confirmed
resident on the GPU after an Ollama restart. This was false: the restart had killed the system
service and launched a user-level server pointing at an empty model directory, which then
crash-looped on a port conflict for roughly two hours; the \texttt{/api/tags} check that appeared
to confirm the models was run \emph{before} the restart, not after. No result in this paper depends
on the local prover having been available during that window, because every Stream~3 goal through
\Cref{sec:towers} was closed by the orchestrating model directly (the orchestrator in this
programme is a model working under the author's direction, not a human checking each step). The error was found only when a
later prover invocation returned \texttt{HTTP~404} for every attempt in four seconds --- an
implausibly fast failure for a model that should take over a minute to cold-load. \textbf{Lesson:}
a status check taken before an action is not evidence about the state after it, and a suspiciously
fast ``success'' is itself a signal to distrust.

\subsection{Identifier ASCII-only}
An early draft named a theorem \texttt{ckn\_$\Lambda$\_le}. This project's own axiom-audit tool
matches declaration names with the regular expression \texttt{[A-Za-z0-9\_'.]+}, which does not
include Greek letters; it silently truncated the name to \texttt{ckn\_} and then reported that no
such declaration existed, rather than reporting a parse failure. No existing declaration anywhere
in this repository's seven prior libraries used a non-ASCII identifier (checked by grep before
this rule was written); this project now keeps that invariant explicitly.

Both incidents, together with the citation error of \Cref{sec:ckn}, are why this paper reports
their discovery rather than only their fixes: the discipline that catches an error is at least as
reportable as the theorem it protects.

%%%%%%%%%%%%%%%%%%%%%%%%%%%%%%%%%%%%%%%%%%%%%%%%%%%%%%%%%%%%%%%%%%%%%%%%%%%%%%%
\section{Open problems for the community}
\label{sec:open}
%%%%%%%%%%%%%%%%%%%%%%%%%%%%%%%%%%%%%%%%%%%%%%%%%%%%%%%%%%%%%%%%%%%%%%%%%%%%%%%
Each item below names the specific gap, not just the headline claim, following this project's
own convention that a precise open problem is worth more than a vague one.

\begin{enumerate}
\item \textbf{The extra-dimension reading is untested at the right precision.}
\Cref{thm:torsion} shows only an order-one margin against a 2020 torsion-balance bound. A
tighter, more recent bound (or the forecast sensitivity of a proposed experiment) formalized the
same way as \Cref{thm:regge} would either exclude or further constrain the reading, resolving the
current ``disfavored, not excluded'' verdict one way or the other.
\item \textbf{The Tier~C identifications are never derived.} $\alpha'=\ell_{\mathrm{micro}}\ell_{\mathrm{macro}}$
and $M=1/\ell_{\mathrm{micro}}$ (\Cref{sec:selfdual}) and $M_1M_2=M_0^2$ (\Cref{sec:towers}) are
this project's modeling choices, stated as hypotheses. A derivation of any of them from an
explicit worldsheet or effective-field-theory argument --- rather than an assumption --- would
upgrade the corresponding results from conditional Tier~A to something closer to Tier~L+A.
\item \textbf{The reduced-vs-non-reduced Planck-mass ambiguity is prose, not Lean.} The factor
$(8\pi)^{1/4}\approx2.2$ separating conventions in \Cref{thm:identity} is discussed but not
formalized as an explicit, named parameter; a Lean statement carrying the convention as an
argument (rather than a fixed numeral) would make every downstream numeric claim
convention-parametric and auditable.
\item \textbf{The SDC itself remains unformalized, for a principled reason (\Cref{sec:towers}).}
An explicit, non-asymptotic toy moduli space --- for instance the circle compactification already
present in Stream~1~\cite{Callens2026Paper7}, where geodesic distance and tower mass are both
concrete --- could support a genuinely non-vacuous, non-asymptotic instance of eq.~(2.100),
where the general statement cannot.
\item \textbf{The defect-network mechanism (\Cref{sec:outofscope}) has no non-vacuous target
identified.} Saikawa's wall-domination time $t_{\mathrm{dom}}=3M_{\mathrm{Pl}}^2/(4A\sigma)$
(eq.~(2.19)) is a genuine arithmetic consequence, not a restatement of the scaling law itself, and
is a candidate for a first Tier~A statement in this area once a specific $\sigma$ is pinned to a
source.
\item \textbf{Numeric brackets are hand-rounded, not interval arithmetic.} Every bracket in
\Cref{sec:selfdual,sec:identity} uses conservatively rounded decimal literals chosen by the
authors. A Lean interval-arithmetic treatment computed directly from the CODATA/Planck-18 constants
would remove the rounding-direction judgment call entirely and tighten every bracket.
\item \textbf{Relevance to $K3\times T^2$ is stated, not derived.} If the ``dark dimension'' reading
of $s$ is taken seriously, the only natural single large direction on this programme's own
compactification is one $T^2$ circle, whose T-dual partner would then be sub-Planckian; nothing in
this repository derives such a decompactification limit or checks it against Stream~2's torus
bound~\cite{Callens2026Paper8}.
\end{enumerate}

%%%%%%%%%%%%%%%%%%%%%%%%%%%%%%%%%%%%%%%%%%%%%%%%%%%%%%%%%%%%%%%%%%%%%%%%%%%%%%%
\section{Proposed research directions}
\label{sec:directions}
%%%%%%%%%%%%%%%%%%%%%%%%%%%%%%%%%%%%%%%%%%%%%%%%%%%%%%%%%%%%%%%%%%%%%%%%%%%%%%%
\begin{enumerate}
\item \textbf{A general ``duality $\iff$ bound-saturation'' lemma.} \Cref{thm:ckniff}'s pattern ---
a monotone bound holds exactly when a length exceeds the self-dual point of a paired duality ---
did not require anything specific to CKN beyond monotonicity in each variable. Formalizing this
as an abstract lemma over an arbitrary duality-bound pair could turn this paper's central
observation into a reusable tool, applicable to other swampland-type UV/IR bounds (the species
scale, the Weak Gravity Conjecture) with the same methodology.
\item \textbf{A systematic library of swampland bounds}, each tested the way \Cref{sec:falsify,sec:readings}
test CKN: state the bound from its source, instantiate it at project-relevant scales, and report
the numeric gap (or lack of one) as a Tier~A fact. This would make ``does hypothesis X survive
bound Y'' a repeatable, auditable procedure rather than a one-off calculation.
\item \textbf{Machine-generated negative controls as a general pattern.} \Cref{prop:negctrl} was
written by hand; a tactic or script that, given a Tier~A theorem with a nontrivial hypothesis,
automatically searches for a counterexample when the hypothesis is dropped would let every future
Stream check its own theorems for the vacuity failure discussed in \Cref{sec:towers} without a
human writing the counterexample each time.
\item \textbf{Confronting the extra-dimension reading with forecast, not only current, data.}
Item~1 of \Cref{sec:open} generalizes: a Lean statement parameterized by an experimental
sensitivity (rather than a fixed 2020 bound) could be re-evaluated as new torsion-balance or
collider results appear, turning \Cref{thm:torsion} into a living, re-checkable claim rather than a
snapshot.
\item \textbf{The compile-time infrastructure recommended for this repository's book
(\texttt{papers/book/chapters/ch38\_open.tex}, items~23--24)} --- attribute-level Tier tracking
and an anti-vacuity linter --- would apply directly to future Stream~3-style work, where every new
theorem in this paper had to be checked by hand against the vacuity failure of \Cref{sec:towers}.
\item \textbf{A first continuous (non-arithmetic-shadow) theorem}, as recommended by an external
review of this programme's book: the classical equivalence of the Nambu--Goto and Polyakov
actions, eliminating the worldsheet metric through its own equation of motion~\cite{Tong2009}
(§1.3, ll.~1255--1280). The pointwise, finite-dimensional step is provable with Mathlib as it
stands; the field-level calculus-of-variations statement is not yet, and closing that gap would be
this programme's first genuinely continuous physics theorem rather than an arithmetic or algebraic
shadow of one.
\end{enumerate}

%%%%%%%%%%%%%%%%%%%%%%%%%%%%%%%%%%%%%%%%%%%%%%%%%%%%%%%%%%%%%%%%%%%%%%%%%%%%%%%
\section{Conclusion}
\label{sec:conclusion}
%%%%%%%%%%%%%%%%%%%%%%%%%%%%%%%%%%%%%%%%%%%%%%%%%%%%%%%%%%%%%%%%%%%%%%%%%%%%%%%
Scale-factor duality and the Cohen--Kaplan--Nelson bound are formalized from pinned sources, and
the dual-tower structure behind the swampland distance conjecture is shown to share scale-factor
duality's algebra rather than adding an independent route. Read literally, the micro/macro
dual-scale hypothesis motivating this stream fails the CKN test by at least $10^{30}$. Read as a
T-dual pair, the two scales meet the CKN bound exactly at their self-dual length, and that length
is --- as an exact identity, not a numerical coincidence --- the observed dark-energy length. It
cannot be the fundamental string's scale, and as an extra-dimension radius it is at present
indistinguishable from an existing, unrelated proposal in the literature. None of this is a proof
or a disproof of the hypothesis; it is a record of what two decades-old pieces of physics actually
imply for it, kernel-checked, with every physical identification labelled for what it is. We
consider the value of this exercise to lie exactly there --- and in the two verification-discipline
failures of \Cref{sec:methods}, each caught by trying to use a claim rather than by reading it, as
much as in the theorems themselves.

%%%%%%%%%%%%%%%%%%%%%%%%%%%%%%%%%%%%%%%%%%%%%%%%%%%%%%%%%%%%%%%%%%%%%%%%%%%%%%%
\appendix
\section{Reproducibility}
\label{app:repro}
%%%%%%%%%%%%%%%%%%%%%%%%%%%%%%%%%%%%%%%%%%%%%%%%%%%%%%%%%%%%%%%%%%%%%%%%%%%%%%%

\subsection{Build}
\begin{center}
\footnotesize\ttfamily
lake build DualScaleStream2 StringTheoryFormalization StringTheoryFoundation \textbackslash\\
  DualScaleM24Formalization DoubleFieldTheory DualScaleValidation Lean5Corpus DualScaleCosmology
\end{center}
3781 jobs, 0 errors (release tag \texttt{v3.4.0}/\texttt{v3.5.x}). The job count is dominated by
Mathlib's own dependency graph, not this library's proof content.

\subsection{Axiom audit}
\begin{center}
\footnotesize\ttfamily
python3 tools/axiom\_audit.py DualScaleCosmology
\end{center}
reports \textbf{31 theorems, 0 failing}, every declaration's axiom footprint limited to
\texttt{propext}, \texttt{Classical.choice}, \texttt{Quot.sound}. Across all eight libraries in
this repository: \textbf{457 theorems, 0 failing} (release \texttt{v3.6.0}, which added one theorem to Stream~2).

\subsection{Statement lock}
\begin{center}
\footnotesize\ttfamily
python3 tools/statement\_lock.py --check \$(find DualScaleCosmology -name '*.lean')
\end{center}
locks all 44 declarations across the 7 files of \texttt{DualScaleCosmology}:
{\small\texttt{ScaleFactorDuality.lean, CKNBound.lean, CKNInstance.lean, DualTower.lean,
SelfDualCutoff.lean, CosmicString.lean, DarkEnergyScale.lean}}.

\subsection{Roadmap coverage}
\texttt{docs/STREAM3\_WORKFLOW.md}~\S6 records every item P3.1--P3.9 as either \checkmark{}
(closed with the theorems reported above) or explicitly closed as out of scope with the reason
stated in \Cref{sec:outofscope}, per the definition of done fixed before the last work package.

\subsection{Sources pinned}
The following are each hash-pinned in \texttt{papers/foundations/MANIFEST.md}, with the line
range quoted at the point of use above:
{\small\texttt{hep-th/9211021, hep-th/9803132, 1903.06239, hep-th/0312067, 1303.5085,
1911.03947, 2002.11761, 2205.12293, 0911.1345, 1703.02576}}.
CODATA~2018 and Planck~2018~VI constants via \texttt{astropy}.

\begin{thebibliography}{99}
\bibitem{Callens2026Paper1} X.~Callens and the SocrateAI Agora Collaboration, ``Dual-Scale Generalized Geometry and Non-Perturbative Moduli Stabilization on $K3\times T^2$,'' \textit{SocrateAI Laboratory for Mathematical Physics \& Formal Verification} (2026).
\bibitem{Callens2026Paper7} X.~Callens and the SocrateAI Agora Collaboration, ``The Dual-Scale String Theory: Mechanized Foundations, Singularity Resolution, Mathieu Moonshine, and a Zero-Free-Parameter Cosmological Conjecture on $K3\times T^2$,'' \textit{SocrateAI Laboratory for Mathematical Physics \& Formal Verification} (2026).
\bibitem{Callens2026Paper8} X.~Callens and the SocrateAI Agora Collaboration, ``Lattices, T-Duality, and Double Field Theory on $K3\times T^2$: A Lean 4 Companion Formalization,'' \textit{SocrateAI Laboratory for Mathematical Physics \& Formal Verification} (2026).
\bibitem{GasperiniVeneziano1993} M.~Gasperini and G.~Veneziano, ``Pre-Big-Bang in String Cosmology,'' \textit{Astropart. Phys.} 1 (1993) 317--339, \textit{arXiv:hep-th/9211021}.
\bibitem{CKN1999} A.~G.~Cohen, D.~B.~Kaplan, and A.~E.~Nelson, ``Effective Field Theory, Black Holes, and the Cosmological Constant,'' \textit{Phys. Rev. Lett.} 82 (1999) 4971--4974, \textit{arXiv:hep-th/9803132}.
\bibitem{PaltiReview2019} E.~Palti, ``The Swampland: Introduction and Review,'' \textit{Fortsch. Phys.} 67 (2019) 1900037, \textit{arXiv:1903.06239}.
\bibitem{CopelandMyersPolchinski2004} E.~J.~Copeland, R.~C.~Myers, and J.~Polchinski, ``Cosmic F- and D-strings,'' \textit{JHEP} 06 (2004) 013, \textit{arXiv:hep-th/0312067}.
\bibitem{Planck2013XXV} Planck Collaboration, ``Planck 2013 results. XXV. Searches for cosmic strings and other topological defects,'' \textit{Astron. Astrophys.} 571 (2014) A25, \textit{arXiv:1303.5085}.
\bibitem{CMS2020} CMS Collaboration, ``Search for high mass dijet resonances with a new background prediction method in proton-proton collisions at $\sqrt{s}=13$~TeV,'' \textit{JHEP} 05 (2020) 033, \textit{arXiv:1911.03947}.
\bibitem{LeeAdelberger2020} J.~G.~Lee, E.~G.~Adelberger, T.~S.~Cook, S.~M.~Fleischer, and B.~R.~Heckel, ``New Test of the Gravitational $1/r^2$ Law at Separations down to $52\,\mu$m,'' \textit{Phys. Rev. Lett.} 124 (2020) 101101, \textit{arXiv:2002.11761}.
\bibitem{MonteroVafaValenzuela2023} M.~Montero, C.~Vafa, and I.~Valenzuela, ``The Dark Dimension and the Swampland,'' \textit{JHEP} 02 (2023) 022, \textit{arXiv:2205.12293}.
\bibitem{CopelandKibble2009} E.~J.~Copeland and T.~W.~B.~Kibble, ``Cosmic Strings and Superstrings,'' \textit{Proc. Roy. Soc. Lond. A} 466 (2010) 623--657, \textit{arXiv:0911.1345}.
\bibitem{Saikawa2017} K.~Saikawa, ``A review of gravitational waves from cosmic domain walls,'' \textit{Universe} 3 (2017) 40, \textit{arXiv:1703.02576}.
\bibitem{Tong2009} D.~Tong, ``Lectures on String Theory,'' \textit{arXiv:0908.0333} (2009).
\bibitem{Planck2018VI} Planck Collaboration, ``Planck 2018 results. VI. Cosmological parameters,'' \textit{Astron. Astrophys.} 641 (2020) A6, \textit{arXiv:1807.06209}.

\end{thebibliography}

\end{document}
